\section{Threats to Validity}\label{sec:threats}

\subsection{Construct Validity}
Our measurement procedures reported false successes four times during the experiment.
Three were caught and fixed before the run; the fourth we found only in a post-submission
audit, and we report it here rather than leave the affected number standing. First, mutation scores of 100\% turned up for suites
that had \emph{failed} on the unmutated code: every mutant was trivially ``killed''
because the suite was already broken. We made a green-on-original check mandatory before
any mutation analysis, so a suite that fails on the baseline source is scored
\textsc{invalid} outright. Second, EvoSuite exited with code 0 (success) even when it
generated no tests at all; a run of 12 functions was logged as a complete success with
zero test files on disk, so the pipeline now checks for the generated file directly rather
than trusting the exit code. Third, JaCoCo silently reported 0\% coverage for EvoSuite
tests, because EvoSuite's separate instrumenting classloader changes class identity
underneath it; disabling \texttt{separateClassLoader} during measurement fixed the
attribution. Fourth, and found too late to change the pipeline, the Java
\texttt{compiled} flag never measured compilation: it was set for every function whose
line range appeared in the JaCoCo report, and JaCoCo enumerates every class in a project
regardless of whether it was loaded. A class that was never touched still yields
\texttt{cb=0}, which the parser records as 0\% coverage with \texttt{compiled=1}.
Recompiling each suite directly --- first with \texttt{javac} against the project
classpath, then confirmed with \texttt{mvn test-compile} on sampled functions from two
repositories --- gives 17/60 rather than 51/60. The failures concentrate on ambiguous
overloads and on targets declared \texttt{private}, of which the sample contains 17. This
does not move any hypothesis test, because RQ1--RQ3 are computed from coverage and
mutation scores rather than from the flag, and the effective count of 5/60 is unaffected;
but it means the Java validity rate reported in Section~\ref{sec:results} should be read
as an upper bound.

One further construct threat: the two Python arms use different scoring instruments ---
\texttt{mutmut} for the \texttt{gpt-4o-mini} suites and our own AST parser for the Pynguin
suites (Section~\ref{sec:method}). Both enforce the same green-on-original gate and score
over the same function ranges, but they define different operator sets, so the Python
RQ2-B comparison really compares scores from two instruments. Fortunately Pynguin's median
coverage on this dataset was 0.0\%, so the direction of that comparison does not hinge on
the instrument choice. This dataset also illustrates why coverage alone is a poor quality
metric: one baseline suite reached 60\% branch coverage on a function while killing none of
its mutants --- which is exactly why we report coverage and mutation score as a pair.

\subsection{Internal Validity}
The pilot run existed to expose integration surprises before the main measurements, and it
was only partly comfortable reading. The initial Java dataset held functions whose file
contents did not match the method named in the metadata --- \texttt{JA-002}'s file held a
\texttt{concat} method, not the \texttt{getOptionValues} the LLM was asked to test, so it
reasoned about the wrong method. Every pilot LLM test was discarded and the dataset was
corrected before the main run. On the Python side, 10/12 pilot measurements were invalid
because their functions had been extracted without their parent modules, an issue
unrelated to the LLM; the later measurements run against whole modules at pinned commits.
Baseline generation had its own defects: EvoSuite's initial 60s search budget against
Randoop's 10s was uncovered, and Randoop's output files were overwritten between runs; both
were fixed before the main measurements, which now use consistent budgets and per-function
output filenames. Several environment and tool mismatches also had to be resolved.
EvoSuite~1.2.0 needs Java~8 for the \texttt{csv}/\texttt{gson}/\texttt{joda} classes while
most code was tested against JDK~11, so the EvoSuite suites carry a mix of JDK~11
(39 functions) and JDK~8 (21 functions) provenance. Multi-release JARs carried Java~21
bytecode that older ASM refused to read (\texttt{Unsupported class file major version 65}),
and \texttt{commons-math} disabled JaCoCo in its own \texttt{pom.xml} via
\texttt{jacoco.skip=true}. Notably, four of the six measurement-pipeline failures reported
exit code 0 --- only the absence of artefacts on disk revealed the problem.

\subsection{External Validity}

\textbf{Single model.} All results are reported for the lightweight
\texttt{gpt-4o-mini-2024-07-18} only and should not be extrapolated beyond it
(amendment~v1.1). \emph{Mitigation applied:} we pinned the exact model snapshot rather than a
floating alias, so the result is attributable to one identifiable system and can be re-run
against it.

\textbf{Corpus skew.} The dataset covers two worlds: Apache-style Java libraries and Python
web libraries (\texttt{flask}, \texttt{requests}). The Python side leans heavily on
\texttt{flask}: 48 of the 60 Python functions come from that package. Functions with heavy
environmental requirements (GUI or network access) are under-represented, and coverage varies
sharply between repositories --- often near-maximal for numeric utility code, near-minimal for
parser- and formatter-style classes. \emph{Mitigation applied:} we report per-repository
provenance for all 120 functions in the replication package, and we read RQ1's 80\% threshold
as repository-relative rather than a universal bar.

\textbf{Single prompt strategy and one complexity band.} The findings describe one one-shot
prompt template at temperature~$=0$ with a single worked exemplar, on functions with
$5 \le CC \le 10$ only. They say nothing about few-shot or chain-of-thought prompting, nor
about functions below or above that band. \emph{Mitigation applied:} the exact prompt and the
band boundaries are fixed and published, so any follow-up varies one factor against a known
reference point rather than re-deriving a baseline.

\textbf{Data contamination.} The ten subject repositories are long-standing open-source
projects (\texttt{flask}, \texttt{requests}, the Apache Commons family, \texttt{gson},
\texttt{jsoup}, \texttt{joda-time}, \texttt{jfreechart}), and their source --- along with
their existing test suites --- was very likely present in the model's training data. Any
apparent competence may therefore be partly recall rather than generation. \emph{Mitigation
applied:} we score every suite by executing it against the project at a pinned commit and by
mutation analysis on the target's own lines, so a memorised-looking test still earns nothing
unless it actually runs and kills mutants; and the measured outcome --- a 0\% all-level median
--- is in any case far too low to be explained by memorisation.

\subsection{Conclusion Validity}

\textbf{Multiple testing.} We run twelve hypothesis tests in total: three per research
question (Python, Java, pooled) for RQ1, RQ2-A and RQ3, and three paired comparisons for
RQ2-B. We report each at the pre-registered $\alpha=0.05$, but a Bonferroni correction for
twelve tests gives $\alpha' = 0.05/12 = 0.0042$, and we state plainly what that does to our
conclusions. \emph{Mitigation applied:} we re-checked every significant result against
$\alpha'$. The two results that go against us survive --- EvoSuite ($p=0.00044$) and Randoop
($p=0.0033$) both still beat the model. The one result that favours the model does
\emph{not}: the win over Pynguin ($p=0.010$) exceeds $\alpha'$ and must be read as
inconclusive once multiplicity is taken into account. No conclusion in this paper rests on
that comparison.

\textbf{Small effective samples.} The all-level analysis uses the full $N=120$, but every
effective-level figure rests on a much smaller subset: $n=23$ (Python coverage), $n=17$
(Python mutation), and $n=5$ (Java coverage). At $n=5$ the Java RQ1 and RQ3 estimates are
indicative only, and the study is underpowered to detect anything but a very large effect
there. \emph{Mitigation applied:} we report the all-level figure beside every effective-level
figure, never the effective one alone, and we state $n$ at every point estimate so no reader
can mistake an $n=5$ median for a population value.

\textbf{Single generation per function.} With one one-shot call per function at
temperature~$=0$, we have no per-function variance estimate and cannot separate model
variability from target difficulty. \emph{Mitigation applied:} temperature was pinned at 0
and the model snapshot pinned to \texttt{gpt-4o-mini-2024-07-18}, so the reported run is at
least deterministic and exactly reproducible from the replication package.

\textbf{Conditional subsets overstate ability.} The effective-level medians
(75\%/50\%/37\%/0\%) are conditioned on suites that both compiled \emph{and} touched the
target, so they overstate raw ability relative to the all-level 0\%. \emph{Mitigation
applied:} both levels are reported at every RQ, and the all-level figure is the one carried
into the conclusions.

With $N=120$ functions (60 per language), we report all measurements with non-parametric
tests where applicable (Wilcoxon, Spearman), at $\alpha=0.05$, and with effect sizes
alongside every $p$-value. The mutation and coverage thresholds were fixed in the approved
protocol before any full run. Every amendment to the pipeline (model settings, dataset,
prompting method, measurement environment) is documented before the measurements it
affects, which is what guards against HARKing. Mutation scores were unavailable (PIT
generated no mutants) for 3/60 of the baseline cells: \texttt{JA-017} (\texttt{jfreechart}),
\texttt{JA-048}, and \texttt{JA-058} (\texttt{jsoup}). These are recorded in
\texttt{metrics\_full.csv} with the mutation-score column empty, since those functions had
no mutants to kill; the coverage figures are unaffected, as they are computed independently
and before the PIT step. The same three functions also appear among the LLM's own
\textsc{invalid}/no-touch Java results (Section~\ref{sec:results}), for an unrelated
reason --- wrong-target invocation. PIT's missing mutants and the LLM's wrong-target
problem were diagnosed independently; the overlap is incidental.

A separate environment issue affected only the Pynguin baseline (Python). Under
Python~3.14 --- the interpreter used everywhere else in the measurement pipeline ---
Pynguin's bytecode instrumentation failed outright, generating empty test suites (0\%
coverage) for every module. Verbose logging traced it to a Pynguin internals problem: its
master--worker process communication broke under Python~3.14, a genuine interpreter-level
incompatibility. The fix was a dedicated Python~3.10 virtual environment for Pynguin
invocations only --- the version most Pynguin releases are validated against. This
adjustment touches none of the pre-registered metrics, thresholds, datasets, or models,
and was applied uniformly across all Pynguin measurements before any data were collected,
which is why it is not an amendment. Even with the fix, Pynguin's median coverage on the
Python dataset stayed at 0.0\% (Section~\ref{sec:results}).

A late audit of that 0.0\% figure shows it blends two distinct causes, and we report the
distinction rather than leave it implicit. For the largest modules Pynguin produced no
output at all --- \texttt{flask/app.py} (1625 LOC) and \texttt{flask/cli.py} (1127 LOC)
among them. We first attributed this to the 90-second budget, since the failure tracks
module size: every module of 385 LOC or fewer yielded a measurable suite, every module of
682 LOC or more yielded none. Re-running with the budget raised to 306\,s refuted that
explanation. Pynguin exits after a few seconds with an empty \texttt{stdout} and
\texttt{rc=2}: its worker process serialises module state with \texttt{dill} and dies on
C-extension type objects that cannot be resolved by name (\texttt{\_ctypes.\_CData} for
these two modules, \texttt{\_json.Scanner} for \texttt{flask/helpers.py}). Module size and
the true cause travel together here because larger Flask modules pull in more such
imports. Disabling the master--worker split lets the search run, so scoring these
functions as 0\% is a statement about tool applicability on this dataset, not about the
quality of the tests Pynguin writes. A second group is more troubling: there, Pynguin did
produce a working suite --- our own harness measures 70.0\% branch coverage on
\texttt{requests.adapters} once the suite is scored per test --- but our
green-on-original check runs against the whole module suite, so one failing test zeroed
every function in that module. Restricted to the 12 functions whose module cleared the
gate, Pynguin's median branch coverage is 40.0\%, not 0.0\%. This is the same
all-or-nothing artefact we describe for the Python LLM arm in Section~\ref{sec:rq4}, and
we had not applied that insight to the baseline. It does not change the RQ1--RQ3
conclusions, which concern \texttt{gpt-4o-mini} rather than Pynguin, but it does mean the
RQ2-B Python comparison should be read as gpt-4o-mini against a baseline that our own
harness partly suppressed. Re-measuring Pynguin with a per-function gate is the first item
of our errata.
